% ----------------------------------------------------------------------
%  2.2.5. Дислокационные структуры
% ----------------------------------------------------------------------
\subsection{Дислокационные структуры}\label{ssec:disl-structures}

\begin{orientir}
Дислокации редко бывают одиночными: они объединяются в группы, и свойства
группы сильно отличаются от суммы свойств отдельных дислокаций. Мы
рассмотрим две противоположные конфигурации: \textbf{скопление}
(неустойчивое, создаёт огромные напряжения; отсюда формула Холла--Петча)
и \textbf{стенку} (устойчивую, почти без дальних полей; из стенок
состоят малоугловые границы). Затем~— как стенки образуют систему
субзёрен (полигонизация) и что такое дисклинации.
\end{orientir}

\begin{glavnoe}
Энергия группы из $n$ дислокаций может заметно отличаться от $n$-кратной
энергии одной дислокации, потому что дальнодействующие поля дислокаций
складываются или гасят друг друга. Чем больше дислокаций, тем важнее эти
\emph{коллективные эффекты}. При малых деформациях (мало дислокаций)
поведение определяют свойства отдельных дислокаций, при больших~—
свойства и взаимодействия больших дислокационных групп.
\end{glavnoe}

\begin{masshtab}[дислокационные структуры]{-7.9/-6.8/{дислокации в стенке}, -6.5/-5/{ячейки и субзёрна}, -5.5/-3.7/{скопления до размера зерна}}
Стенки (субграницы) состоят из дислокаций, отстоящих на $h=b/\varphi$:
десятки--сотни нанометров при углах $0{,}1$--$1^\circ$. Ячейки и субзёрна,
которые ими ограничены,~— 0,3--10~мкм, т.\,е. в одном зерне их сотни и
тысячи. Длина плоского скопления ограничена размером зерна
(1--100~мкм)~— отсюда зависимость предела текучести от размера зерна
(формула Холла--Петча).
\end{masshtab}

\subsubsection{Скопление дислокаций}

\begin{opredelenie}
\term{Скопление дислокаций}~— группа одноимённых дислокаций, скользящих
в одной плоскости и остановленных препятствием (границей зерна,
межфазной границей и т.\,п.).
\end{opredelenie}

\begin{figure}[ht]
\centering
\begin{tikzpicture}[scale=0.9]
  \fill[fmGray!25] (-0.5,-1.0) rectangle (0,1.0);
  \draw[fmGray, line width=1pt] (0,-1.0) -- (0,1.0);
  \node[slbl, align=center] at (-0.25,-1.4) {препятствие};
  \draw[fmGray, dashed] (0,0) -- (11.2,0);
  \foreach \i in {1,...,9}{
    \pgfmathsetmacro{\xx}{0.25+10.5*((\i-0.6)/9)^2}
    \dislT[fmBlue]{(\xx,0)}{0}}
  \draw[arr, fmRed] (12.0,0.45) -- (10.9,0.45) node[left, slbl, pos=0]{$\tau$\ \ };
  \draw[dimline] (0,-0.6) -- node[below, slbl]{$L$} (10.75,-0.6);
  \node[slbl, above] at (0.26,0.35) {1};
  \node[slbl, above] at (0.73,0.35) {2};
  \node[slbl, above] at (10.73,0.35) {$N$};
\end{tikzpicture}
\caption{Скопление $N$ краевых дислокаций у препятствия. Внешнее
напряжение $\tau$ прижимает дислокации к препятствию, а взаимное
отталкивание их раздвигает; у головы скопления они расположены гуще всего}
\label{fig:pileup}
\end{figure}

\paragraph{Скопление энергетически невыгодно}
Издалека ($R\gg L$, где $L$~— длина скопления) скопление выглядит как
одна «супердислокация» с вектором Бюргерса $Nb$. Её энергия
$\propto(Nb)^2=N^2b^2$, тогда как энергия $N$ разнесённых дислокаций
$\propto Nb^2$. Скопление~— неустойчивая конфигурация, которая стремится
распасться. Поэтому оно не может возникнуть за счёт перегруппировки
дислокаций (любая перегруппировка уменьшает энергию). Скопление~—
результат работы источника, испускающего дислокации в одной плоскости
навстречу препятствию.

При повышенных температурах дислокации скопления переползают в
соседние плоскости скольжения и выстраиваются в \emph{стенку}
(рис.~\ref{fig:pileup-to-wall}). Переползание требует диффузии, поэтому
скопления живут долго лишь при низких температурах.

\begin{figure}[ht]
\centering
\begin{tikzpicture}[scale=0.8]
  \draw[fmGray, dashed] (0,0) -- (6,0);
  \foreach \i in {1,...,5}{\pgfmathsetmacro{\xx}{0.4+5*((\i-0.6)/5)^2}\dislT[fmBlue]{(\xx,0)}{0}}
  \draw[fmGray, line width=1pt] (0,-1.5) -- (0,1.5);
  \draw[arr, fmGray, line width=1pt] (6.6,0) -- (7.6,0);
  \begin{scope}[shift={(9,0)}]
    \foreach \y in {-2,...,2}{\draw[fmGray, dashed] (-0.8,\y*0.65) -- (0.8,\y*0.65); \dislT[fmBlue]{(0,\y*0.65)}{0}}
  \end{scope}
\end{tikzpicture}
\caption{Перестройка скопления дислокаций в стенку путём переползания
(при достаточно высокой температуре)}
\label{fig:pileup-to-wall}
\end{figure}

\paragraph{Скопление~— концентратор напряжений}
Пусть все $N$ дислокаций скопления заняли равновесные положения. На
каждую действует внешняя сила $\tau b$ (на единицу длины). Последнюю
дислокацию сила $\tau b$ толкает к препятствию, но она стоит~— значит,
остальные дислокации отталкивают её с силой $\tau b$. По третьему закону
Ньютона она сама давит на остальных с добавочной силой $\tau b$.
Рассуждая так дальше, видим, что $m$ «хвостовых» дислокаций давят на
остальные с силой $m\tau b$. На головную дислокацию действуют внешняя
сила и давление всех $N-1$ остальных:
\begin{keyeq}\label{eq:pileup-head}
F_1=\tau b+(N-1)\tau b=N\tau b .
\end{keyeq}
\begin{glavnoe}
Скопление из $N$ дислокаций концентрирует приложенное напряжение в $N$
раз. С такой же силой голова скопления давит на препятствие. Это может
вызвать прорыв препятствия, срабатывание источника в соседнем зерне или
зарождение микротрещины.
\end{glavnoe}

Ситуация напоминает цепочку из $N$ грузов, связанных нитями: на последней
нити висит вес $NP$. Разница в том, что грузы связаны только с соседями,
а дислокации~— все со всеми (поля напряжений дальнодействующие).
Поэтому точно найти положения всех дислокаций непросто, но оценки
получаются легко.

Обозначим $A=G/[2\pi(1-\nu)]$; сила отталкивания двух одноимённых краевых
дислокаций в одной плоскости на расстоянии $x$ равна $Ab^2/x$
(п.~\ref{ssec:disl-elastic}).

\paragraph{Расстояние между головными дислокациями}
На вторую дислокацию к препятствию действуют внешняя сила $\tau b$ и
давление остальных $N-2$ дислокаций; поскольку расстояние $d$ между двумя
головными мало по сравнению с $L$, считаем, что это давление делится
между первой и второй дислокациями поровну: $\frac{N-2}{2}\tau b$. В
обратную сторону её отталкивает первая дислокация с силой $Ab^2/d$.
Равновесие:
\[
\frac{Ab^2}{d}=\tau b+\frac{N-2}{2}\tau b=\frac N2\tau b,\qquad
d=\frac{2Ab}{N\tau} .
\]

\begin{primer}[Скопление в алюминии]
$A\approx600$~кгс/мм$^2$, $N=25$, $\tau=6$~кгс/мм$^2$:
$d/b=2\cdot600/(25\cdot6)=8$. Головные дислокации отстоят друг от друга
всего на несколько межатомных расстояний.
\end{primer}

\paragraph{Длина скопления}
На последнюю дислокацию к препятствию действует только внешняя сила $\tau
b$, а обратно~— отталкивание всех остальных. Они далеко, поэтому их можно
заменить одной супердислокацией $(N-1)b$ в середине скопления (на
расстоянии $L/2$):
\begin{equation}\label{eq:pileup-L}
\frac{A(N-1)b^2}{L/2}=\tau b\quad\Longrightarrow\quad L\approx\frac{2NAb}{\tau} .
\end{equation}

\subsubsection{Формула Холла--Петча}

В поликристалле длина скопления не может превышать размер зерна $l$.
Из~\eqref{eq:pileup-L} при заданном напряжении $\tau$ на длине $l$
помещается $N=l\tau/(2Ab)$ дислокаций, и у границы зерна возникает
напряжение
\[
\tau_1=N\tau=\frac{l\tau^2}{2Ab} .
\]
Пусть для какого-то процесса у границы (прорыва границы, запуска
источника в соседнем зерне) нужно напряжение $\tau_1$. Тогда приложенное
напряжение должно быть не меньше $\tau=\sqrt{2Ab\tau_1/l}$. Добавим
напряжение трения $\tau_0$, нужное просто для движения дислокации по
кристаллу (рельеф Пайерлса, примеси, другие дефекты), и обозначим
$k=\sqrt{2Ab\tau_1}$\ispr{В оригинале (с.~269) записано «$2D\tau_1=k$»;
из вывода следует $k=\sqrt{2Ab\tau_1}$ (корень и множитель $b$
потеряны).}:
\begin{keyeq}\label{eq:hall-petch}
\tau_{\text{т}}=\tau_0+k\,l^{-1/2}\qquad\text{(формула Холла--Петча)} .
\end{keyeq}

\begin{glavnoe}
Предел текучести растёт при измельчении зерна: $\sigma_\text{т}\propto
l^{-1/2}$. Граница зерна~— барьер для скоплений, и чем мельче зерно, тем
короче скопления и тем меньше концентрация напряжений у границ.
Измельчение зерна~— один из главных способов упрочнения металлов
(который, в отличие от многих других, не снижает, а повышает вязкость).
\end{glavnoe}

\begin{primer}[Упрочнение измельчением зерна]
Для мягкой стали типичны значения $\sigma_0\approx100$~МПа,
$k_y\approx0{,}6$~МПа$\cdot$м$^{1/2}$ (формула Холла--Петча для
нормального напряжения). При размере зерна $l=100$~мкм:
$\sigma_\text{т}=100+0{,}6/\sqrt{10^{-4}}=160$~МПа; при $l=1$~мкм:
$\sigma_\text{т}=100+0{,}6/\sqrt{10^{-6}}=700$~МПа. Уменьшение зерна в
100 раз повышает предел текучести более чем вчетверо.
\end{primer}

Препятствиями для скоплений служат не только границы зёрен, но и стенки
дислокаций; тогда в формулу~\eqref{eq:hall-petch} вместо размера зерна
входит размер субзерна.

\begin{podrobnee}[Коллективные эффекты и континуальное описание скоплений]
\emph{Подвижность скопления.} Чтобы задержать скопление, нужно очень
мощное препятствие. Если скопление движется через хаотично разбросанные
препятствия, его дислокации «помогают» друг другу, концентрируя напряжения
у самых сильных препятствий. Опыт показывает, что повышение напряжения
всего в 3--4 раза увеличивает скорость дислокаций на несколько порядков,
поэтому скопления подвижнее одиночных дислокаций. Особенно это заметно
при переменной нагрузке: при разгрузке дислокации отходят от части
препятствий, а при новом нагружении у оставшихся возникает повышенная
концентрация напряжений, и они преодолеваются «динамически»~—
последовательно, а не параллельно, что из-за нелинейной зависимости
$v(\tau)$ эффективнее.

\emph{Континуальное приближение.} Чтобы найти поле скопления, удобно
«размазать» дислокации: каждую разбить на бесконечно много бесконечно
малых дислокаций с линейной плотностью $n(x)$ (число дислокаций на
единицу длины), причём $\int_0^Ln(x)\,\d x=N$. Из условия равновесия
каждой бесконечно малой дислокации для заторможенного скопления в поле
постоянного напряжения получается
\[
n(x)=\frac{\tau}{\pi Ab}\sqrt{\frac{L-x}{x}} ,
\]
и интегрирование даёт снова $L=2NAb/\tau$. У головы скопления
($x\to0$) плотность дислокаций неограниченно растёт. Зная $n(x)$, можно
найти напряжения в любой точке вне скопления, например
\[
\sigma_{xx}(x,y)=-Ab\int_0^Ln(x')\,\frac{y\bigl[3(x-x')^2+y^2\bigr]}{\bigl[(x-x')^2+y^2\bigr]^2}\,\d x' .
\]
\end{podrobnee}

\subsubsection{Стенка дислокаций}

\begin{opredelenie}
\term{Дислокационная стенка}~— ряд одноимённых краевых дислокаций,
расположенных одна над другой в параллельных плоскостях скольжения.
\end{opredelenie}

Стенка, в отличие от скопления, \emph{устойчива}
(п.~\ref{ssec:disl-elastic}, рис.~\ref{fig:edge-edge}). Её энергия меньше
энергии того же числа отдельных дислокаций: поле сжатия каждой дислокации
накладывается на поле растяжения дислокаций, расположенных выше, а поле
растяжения~— на поле сжатия расположенных ниже. В результате поля
практически гасятся на расстояниях порядка расстояния $h$ между
дислокациями, и радиусом обрезания служит $h$:
\begin{equation}\label{eq:wall-energy}
E_{\text{ст}}\approx n\,\frac{Gb^2}{4\pi(1-\nu)}\ln\frac{h}{r_0}
\end{equation}
(без учёта энергии ядер). Энергия дислокаций в стенке в 3--4 раза
меньше, чем у хаотически распределённых.

\paragraph{Свойства стенки}
\begin{itemize}
\item Напряжения отдельных дислокаций гасятся на расстояниях больше $h$.
Стенка~— плоский (поверхностный) дефект толщиной $\sim2h$.
\item Каждая дислокация стенки лежит в своей плоскости скольжения.
Чтобы выстроиться в стенку, дислокации, испущенные хаотически или в одной
плоскости, должны \emph{переползти}. Поэтому стенки~— типичный элемент
структуры только при достаточно высоких температурах:
$T\gtrsim(0{,}3\text{--}0{,}5)\,\Tpl$.
\end{itemize}

\paragraph{Стенка~— граница наклона}
Краевая дислокация~— край лишней полуплоскости. Стенка~— это пачка
полуплоскостей, оканчивающихся на разных уровнях. На расстояниях больше
$h$ от стенки их действие эквивалентно вставленному клину: $n$
полуплоскостей толщиной $b$ на высоте $nh$. Поэтому части кристалла по
обе стороны стенки повёрнуты друг относительно друга на угол
(рис.~\ref{fig:tilt-wall})
\begin{keyeq}\label{eq:tilt-angle}
\varphi\approx\frac bh .
\end{keyeq}

\begin{figure}[ht]
\centering
\begin{tikzpicture}[scale=0.55]
  % Сгенерировано tools/gen_tilt_boundary.py — не редактировать вручную
\begin{scope}
\clip (-6.30,-6.30) rectangle (6.30,6.30);
\draw[atomLine!55, line width=0.45pt] (-7.059,-0.102)--(-7.178,0.897) (-7.059,-0.102)--(-6.940,-1.008) (-7.059,-0.102)--(-6.060,-0.008) (-7.178,0.897)--(-7.297,1.896) (-7.178,0.897)--(-6.180,0.992) (-7.297,1.896)--(-6.299,1.990) (-6.535,3.893)--(-6.653,4.890) (-6.535,3.893)--(-6.419,2.989) (-6.535,3.893)--(-5.539,3.987) (-6.653,4.890)--(-6.772,5.887) (-6.653,4.890)--(-5.658,4.984) (-6.772,5.887)--(-6.890,6.884) (-6.772,5.887)--(-5.778,5.980) (-6.890,6.884)--(-5.897,6.976) (-7.222,-7.020)--(-6.216,-6.922) (-6.581,-4.011)--(-6.701,-3.010) (-6.581,-4.011)--(-6.457,-4.916) (-6.581,-4.011)--(-5.578,-3.913) (-6.701,-3.010)--(-6.820,-2.008) (-6.701,-3.010)--(-5.699,-2.911) (-6.820,-2.008)--(-6.940,-1.008) (-6.820,-2.008)--(-5.819,-1.910) (-6.940,-1.008)--(-5.940,-0.910) (-6.060,-0.008)--(-6.180,0.992) (-6.060,-0.008)--(-5.940,-0.910) (-6.060,-0.008)--(-5.060,0.090) (-6.180,0.992)--(-6.299,1.990) (-6.180,0.992)--(-5.181,1.090) (-6.299,1.990)--(-6.419,2.989) (-6.299,1.990)--(-5.301,2.089) (-6.419,2.989)--(-5.422,3.087) (-5.539,3.987)--(-5.658,4.984) (-5.539,3.987)--(-5.422,3.087) (-5.539,3.987)--(-4.543,4.084) (-5.658,4.984)--(-5.778,5.980) (-5.658,4.984)--(-4.663,5.081) (-5.778,5.980)--(-5.897,6.976) (-5.778,5.980)--(-4.784,6.078) (-5.897,6.976)--(-4.904,7.073) (-6.216,-6.922)--(-6.337,-5.919) (-6.216,-6.922)--(-5.211,-6.820) (-6.337,-5.919)--(-6.457,-4.916) (-6.337,-5.919)--(-5.331,-5.816) (-6.457,-4.916)--(-5.453,-4.814) (-5.578,-3.913)--(-5.699,-2.911) (-5.578,-3.913)--(-5.453,-4.814) (-5.578,-3.913)--(-4.575,-3.811) (-5.699,-2.911)--(-5.819,-1.910) (-5.699,-2.911)--(-4.697,-2.808) (-5.819,-1.910)--(-5.940,-0.910) (-5.819,-1.910)--(-4.817,-1.807) (-5.940,-0.910)--(-4.939,-0.807) (-5.060,0.090)--(-5.181,1.090) (-5.060,0.090)--(-4.939,-0.807) (-5.060,0.090)--(-4.061,0.193) (-5.181,1.090)--(-5.301,2.089) (-5.181,1.090)--(-4.183,1.193) (-5.301,2.089)--(-5.422,3.087) (-5.301,2.089)--(-4.303,2.192) (-5.422,3.087)--(-4.425,3.189) (-4.543,4.084)--(-4.663,5.081) (-4.543,4.084)--(-4.425,3.189) (-4.543,4.084)--(-3.547,4.186) (-4.663,5.081)--(-4.784,6.078) (-4.663,5.081)--(-3.669,5.184) (-4.784,6.078)--(-4.904,7.073) (-4.784,6.078)--(-3.789,6.180) (-4.904,7.073)--(-3.910,7.174) (-5.211,-6.820)--(-5.331,-5.816) (-5.211,-6.820)--(-4.208,-6.713) (-5.331,-5.816)--(-5.453,-4.814) (-5.331,-5.816)--(-4.325,-5.709) (-5.453,-4.814)--(-4.447,-4.708) (-4.575,-3.811)--(-4.697,-2.808) (-4.575,-3.811)--(-4.447,-4.708) (-4.575,-3.811)--(-3.574,-3.706) (-4.697,-2.808)--(-4.817,-1.807) (-4.697,-2.808)--(-3.697,-2.700) (-4.817,-1.807)--(-4.939,-0.807) (-4.817,-1.807)--(-3.814,-1.699) (-4.939,-0.807)--(-3.937,-0.701) (-4.061,0.193)--(-4.183,1.193) (-4.061,0.193)--(-3.937,-0.701) (-4.061,0.193)--(-3.063,0.299) (-4.183,1.193)--(-4.303,2.192) (-4.183,1.193)--(-3.186,1.301) (-4.303,2.192)--(-4.425,3.189) (-4.303,2.192)--(-3.303,2.300) (-4.425,3.189)--(-3.426,3.294) (-3.547,4.186)--(-3.669,5.184) (-3.547,4.186)--(-3.426,3.294) (-3.547,4.186)--(-2.553,4.292) (-3.669,5.184)--(-3.789,6.180) (-3.669,5.184)--(-2.675,5.291) (-3.789,6.180)--(-3.910,7.174) (-3.789,6.180)--(-2.792,6.287) (-3.910,7.174)--(-2.915,7.279) (-4.208,-6.713)--(-4.325,-5.709) (-4.208,-6.713)--(-3.208,-6.598) (-4.325,-5.709)--(-4.447,-4.708) (-4.325,-5.709)--(-3.315,-5.594) (-4.447,-4.708)--(-3.438,-4.601) (-3.574,-3.706)--(-3.697,-2.700) (-3.574,-3.706)--(-3.438,-4.601) (-3.574,-3.706)--(-2.578,-3.598) (-3.697,-2.700)--(-3.814,-1.699) (-3.697,-2.700)--(-2.701,-2.585) (-3.814,-1.699)--(-3.937,-0.701) (-3.814,-1.699)--(-2.807,-1.583) (-3.937,-0.701)--(-2.930,-0.594) (-3.063,0.299)--(-3.186,1.301) (-3.063,0.299)--(-2.930,-0.594) (-3.063,0.299)--(-2.070,0.406) (-3.186,1.301)--(-3.303,2.300) (-3.186,1.301)--(-2.193,1.417) (-3.303,2.300)--(-3.426,3.294) (-3.303,2.300)--(-2.299,2.415) (-3.426,3.294)--(-2.422,3.402) (-2.553,4.292)--(-2.675,5.291) (-2.553,4.292)--(-2.422,3.402) (-2.553,4.292)--(-1.562,4.399) (-2.675,5.291)--(-2.792,6.287) (-2.675,5.291)--(-1.685,5.406) (-2.792,6.287)--(-2.915,7.279) (-2.792,6.287)--(-1.792,6.402) (-3.208,-6.598)--(-3.315,-5.594) (-3.208,-6.598)--(-2.221,-6.473) (-3.315,-5.594)--(-3.438,-4.601) (-3.315,-5.594)--(-2.293,-5.469) (-3.438,-4.601)--(-2.411,-4.495) (-2.578,-3.598)--(-2.701,-2.585) (-2.578,-3.598)--(-2.411,-4.495) (-2.578,-3.598)--(-1.598,-3.492) (-2.701,-2.585)--(-2.807,-1.583) (-2.701,-2.585)--(-1.716,-2.460) (-2.807,-1.583)--(-2.930,-0.594) (-2.807,-1.583)--(-1.789,-1.458) (-2.930,-0.594)--(-1.906,-0.487) (-2.070,0.406)--(-2.193,1.417) (-2.070,0.406)--(-1.906,-0.487) (-2.070,0.406)--(-1.094,0.513) (-2.193,1.417)--(-2.299,2.415) (-2.193,1.417)--(-1.211,1.542) (-2.299,2.415)--(-2.422,3.402) (-2.299,2.415)--(-1.284,2.540) (-2.422,3.402)--(-1.402,3.508) (-1.562,4.399)--(-1.685,5.406) (-1.562,4.399)--(-1.402,3.508) (-1.562,4.399)--(-0.589,4.505) (-1.685,5.406)--(-1.792,6.402) (-1.685,5.406)--(-0.707,5.531) (-1.792,6.402)--(-0.779,6.527) (-2.221,-6.473)--(-2.293,-5.469) (-2.221,-6.473)--(-1.244,-6.356) (-2.293,-5.469)--(-2.411,-4.495) (-2.293,-5.469)--(-1.261,-5.352) (-2.411,-4.495)--(-1.317,-4.365) (-1.598,-3.492)--(-1.716,-2.460) (-1.598,-3.492)--(-1.317,-4.365) (-1.598,-3.492)--(-0.686,-3.362) (-1.716,-2.460)--(-1.789,-1.458) (-1.716,-2.460)--(-0.742,-2.343) (-1.789,-1.458)--(-1.906,-0.487) (-1.789,-1.458)--(-0.759,-1.341) (-1.906,-0.487)--(-0.816,-0.357) (-1.094,0.513)--(-1.211,1.542) (-1.094,0.513)--(-0.816,-0.357) (-1.094,0.513)--(-0.184,0.643) (-1.211,1.542)--(-1.284,2.540) (-1.211,1.542)--(-0.241,1.659) (-1.284,2.540)--(-1.402,3.508) (-1.284,2.540)--(-0.258,2.657) (-1.402,3.508)--(-0.314,3.638) (-0.589,4.505)--(-0.707,5.531) (-0.589,4.505)--(-0.314,3.638) (-0.589,4.505)--(0.317,4.635) (-0.707,5.531)--(-0.779,6.527) (-0.707,5.531)--(0.261,5.648) (-0.779,6.527)--(0.244,6.644) (-1.244,-6.356)--(-1.261,-5.352) (-1.244,-6.356)--(-0.256,-6.356) (-1.261,-5.352)--(-1.317,-4.365) (-1.261,-5.352)--(-0.239,-5.352) (-1.317,-4.365)--(-0.686,-3.362) (-1.317,-4.365)--(-0.183,-4.365) (-0.686,-3.362)--(-0.742,-2.343) (-0.686,-3.362)--(-0.183,-4.365) (-0.686,-3.362)--(0.186,-3.362) (-0.742,-2.343)--(-0.759,-1.341) (-0.742,-2.343)--(0.242,-2.343) (-0.759,-1.341)--(-0.816,-0.357) (-0.759,-1.341)--(0.259,-1.341) (-0.816,-0.357)--(-0.184,0.643) (-0.816,-0.357)--(0.316,-0.357) (-0.184,0.643)--(-0.241,1.659) (-0.184,0.643)--(0.316,-0.357) (-0.184,0.643)--(0.684,0.643) (-0.241,1.659)--(-0.258,2.657) (-0.241,1.659)--(0.741,1.659) (-0.258,2.657)--(-0.314,3.638) (-0.258,2.657)--(0.758,2.657) (-0.314,3.638)--(0.317,4.635) (-0.314,3.638)--(0.814,3.638) (0.317,4.635)--(0.261,5.648) (0.317,4.635)--(0.814,3.638) (0.317,4.635)--(1.183,4.635) (0.261,5.648)--(0.244,6.644) (0.261,5.648)--(1.239,5.648) (0.244,6.644)--(1.256,6.644) (-0.256,-6.356)--(-0.239,-5.352) (-0.256,-6.356)--(0.721,-6.473) (-0.239,-5.352)--(-0.183,-4.365) (-0.239,-5.352)--(0.793,-5.469) (-0.183,-4.365)--(0.186,-3.362) (-0.183,-4.365)--(0.911,-4.495) (0.186,-3.362)--(0.242,-2.343) (0.186,-3.362)--(1.098,-3.492) (0.242,-2.343)--(0.259,-1.341) (0.242,-2.343)--(1.216,-2.460) (0.259,-1.341)--(0.316,-0.357) (0.259,-1.341)--(1.289,-1.458) (0.316,-0.357)--(0.684,0.643) (0.316,-0.357)--(1.406,-0.487) (0.684,0.643)--(0.741,1.659) (0.684,0.643)--(1.594,0.513) (0.741,1.659)--(0.758,2.657) (0.741,1.659)--(1.711,1.542) (0.758,2.657)--(0.814,3.638) (0.758,2.657)--(1.784,2.540) (0.814,3.638)--(1.183,4.635) (0.814,3.638)--(1.902,3.508) (1.183,4.635)--(1.239,5.648) (1.183,4.635)--(2.089,4.505) (1.239,5.648)--(1.256,6.644) (1.239,5.648)--(2.207,5.531) (1.256,6.644)--(2.279,6.527) (0.721,-6.473)--(0.793,-5.469) (0.721,-6.473)--(1.708,-6.598) (0.793,-5.469)--(0.911,-4.495) (0.793,-5.469)--(1.815,-5.594) (0.911,-4.495)--(1.098,-3.492) (0.911,-4.495)--(1.938,-4.601) (1.098,-3.492)--(1.216,-2.460) (1.098,-3.492)--(2.078,-3.598) (1.216,-2.460)--(1.289,-1.458) (1.216,-2.460)--(2.201,-2.585) (1.289,-1.458)--(1.406,-0.487) (1.289,-1.458)--(2.307,-1.583) (1.406,-0.487)--(1.594,0.513) (1.406,-0.487)--(2.430,-0.594) (1.594,0.513)--(1.711,1.542) (1.594,0.513)--(2.570,0.406) (1.711,1.542)--(1.784,2.540) (1.711,1.542)--(2.693,1.417) (1.784,2.540)--(1.902,3.508) (1.784,2.540)--(2.799,2.415) (1.902,3.508)--(2.089,4.505) (1.902,3.508)--(2.922,3.402) (2.089,4.505)--(2.207,5.531) (2.089,4.505)--(3.062,4.399) (2.207,5.531)--(2.279,6.527) (2.207,5.531)--(3.185,5.406) (2.279,6.527)--(3.292,6.402) (1.708,-6.598)--(1.815,-5.594) (1.708,-6.598)--(2.708,-6.713) (1.815,-5.594)--(1.938,-4.601) (1.815,-5.594)--(2.825,-5.709) (1.938,-4.601)--(2.078,-3.598) (1.938,-4.601)--(2.947,-4.708) (2.078,-3.598)--(2.201,-2.585) (2.078,-3.598)--(3.074,-3.706) (2.201,-2.585)--(2.307,-1.583) (2.201,-2.585)--(3.197,-2.700) (2.307,-1.583)--(2.430,-0.594) (2.307,-1.583)--(3.314,-1.699) (2.430,-0.594)--(2.570,0.406) (2.430,-0.594)--(3.437,-0.701) (2.570,0.406)--(2.693,1.417) (2.570,0.406)--(3.563,0.299) (2.693,1.417)--(2.799,2.415) (2.693,1.417)--(3.686,1.301) (2.799,2.415)--(2.922,3.402) (2.799,2.415)--(3.803,2.300) (2.922,3.402)--(3.062,4.399) (2.922,3.402)--(3.926,3.294) (3.062,4.399)--(3.185,5.406) (3.062,4.399)--(4.053,4.292) (3.185,5.406)--(3.292,6.402) (3.185,5.406)--(4.175,5.291) (3.292,6.402)--(4.292,6.287) (2.708,-6.713)--(2.825,-5.709) (2.708,-6.713)--(3.711,-6.820) (2.825,-5.709)--(2.947,-4.708) (2.825,-5.709)--(3.831,-5.816) (2.947,-4.708)--(3.074,-3.706) (2.947,-4.708)--(3.953,-4.814) (3.074,-3.706)--(3.197,-2.700) (3.074,-3.706)--(4.075,-3.811) (3.197,-2.700)--(3.314,-1.699) (3.197,-2.700)--(4.197,-2.808) (3.314,-1.699)--(3.437,-0.701) (3.314,-1.699)--(4.317,-1.807) (3.437,-0.701)--(3.563,0.299) (3.437,-0.701)--(4.439,-0.807) (3.563,0.299)--(3.686,1.301) (3.563,0.299)--(4.561,0.193) (3.686,1.301)--(3.803,2.300) (3.686,1.301)--(4.683,1.193) (3.803,2.300)--(3.926,3.294) (3.803,2.300)--(4.803,2.192) (3.926,3.294)--(4.053,4.292) (3.926,3.294)--(4.925,3.189) (4.053,4.292)--(4.175,5.291) (4.053,4.292)--(5.047,4.186) (4.175,5.291)--(4.292,6.287) (4.175,5.291)--(5.169,5.184) (4.292,6.287)--(4.415,7.279) (4.292,6.287)--(5.289,6.180) (4.415,7.279)--(5.410,7.174) (3.711,-6.820)--(3.831,-5.816) (3.711,-6.820)--(4.716,-6.922) (3.831,-5.816)--(3.953,-4.814) (3.831,-5.816)--(4.837,-5.919) (3.953,-4.814)--(4.075,-3.811) (3.953,-4.814)--(4.957,-4.916) (4.075,-3.811)--(4.197,-2.808) (4.075,-3.811)--(5.078,-3.913) (4.197,-2.808)--(4.317,-1.807) (4.197,-2.808)--(5.199,-2.911) (4.317,-1.807)--(4.439,-0.807) (4.317,-1.807)--(5.319,-1.910) (4.439,-0.807)--(4.561,0.193) (4.439,-0.807)--(5.440,-0.910) (4.561,0.193)--(4.683,1.193) (4.561,0.193)--(5.560,0.090) (4.683,1.193)--(4.803,2.192) (4.683,1.193)--(5.681,1.090) (4.803,2.192)--(4.925,3.189) (4.803,2.192)--(5.801,2.089) (4.925,3.189)--(5.047,4.186) (4.925,3.189)--(5.922,3.087) (5.047,4.186)--(5.169,5.184) (5.047,4.186)--(6.043,4.084) (5.169,5.184)--(5.289,6.180) (5.169,5.184)--(6.163,5.081) (5.289,6.180)--(5.410,7.174) (5.289,6.180)--(6.284,6.078) (5.410,7.174)--(6.404,7.073) (4.716,-6.922)--(4.837,-5.919) (4.716,-6.922)--(5.722,-7.020) (4.837,-5.919)--(4.957,-4.916) (4.837,-5.919)--(5.842,-6.016) (4.957,-4.916)--(5.078,-3.913) (4.957,-4.916)--(5.961,-5.013) (5.078,-3.913)--(5.199,-2.911) (5.078,-3.913)--(6.081,-4.011) (5.199,-2.911)--(5.319,-1.910) (5.199,-2.911)--(6.201,-3.010) (5.319,-1.910)--(5.440,-0.910) (5.319,-1.910)--(6.320,-2.008) (5.440,-0.910)--(5.560,0.090) (5.440,-0.910)--(6.440,-1.008) (5.560,0.090)--(5.681,1.090) (5.560,0.090)--(6.560,-0.008) (5.681,1.090)--(5.801,2.089) (5.681,1.090)--(6.680,0.992) (5.801,2.089)--(5.922,3.087) (5.801,2.089)--(6.799,1.990) (5.922,3.087)--(6.043,4.084) (5.922,3.087)--(6.919,2.989) (6.043,4.084)--(6.163,5.081) (6.043,4.084)--(7.039,3.987) (6.163,5.081)--(6.284,6.078) (6.163,5.081)--(7.158,4.984) (6.284,6.078)--(6.404,7.073) (6.284,6.078)--(7.278,5.980) (5.722,-7.020)--(5.842,-6.016) (5.722,-7.020)--(6.728,-7.113) (5.842,-6.016)--(5.961,-5.013) (5.842,-6.016)--(6.847,-6.110) (5.961,-5.013)--(6.081,-4.011) (5.961,-5.013)--(6.965,-5.107) (6.081,-4.011)--(6.201,-3.010) (6.081,-4.011)--(7.084,-4.105) (6.201,-3.010)--(6.320,-2.008) (6.201,-3.010)--(7.203,-3.104) (6.320,-2.008)--(6.440,-1.008) (6.440,-1.008)--(6.560,-0.008) (6.560,-0.008)--(6.680,0.992) (6.680,0.992)--(6.799,1.990) (6.799,1.990)--(6.919,2.989) (6.919,2.989)--(7.039,3.987) (7.039,3.987)--(7.158,4.984) (7.158,4.984)--(7.278,5.980) (6.728,-7.113)--(6.847,-6.110) (6.847,-6.110)--(6.965,-5.107) (6.965,-5.107)--(7.084,-4.105) (7.084,-4.105)--(7.203,-3.104);
\fill[atomA, draw=atomLine, line width=0.35pt] (-7.059,-0.102) circle (0.20);
\fill[atomA, draw=atomLine, line width=0.35pt] (-7.178,0.897) circle (0.20);
\fill[atomA, draw=atomLine, line width=0.35pt] (-7.297,1.896) circle (0.20);
\fill[atomA, draw=atomLine, line width=0.35pt] (-6.535,3.893) circle (0.20);
\fill[atomA, draw=atomLine, line width=0.35pt] (-6.653,4.890) circle (0.20);
\fill[atomA, draw=atomLine, line width=0.35pt] (-6.772,5.887) circle (0.20);
\fill[atomA, draw=atomLine, line width=0.35pt] (-6.890,6.884) circle (0.20);
\fill[atomA, draw=atomLine, line width=0.35pt] (-7.222,-7.020) circle (0.20);
\fill[atomA, draw=atomLine, line width=0.35pt] (-6.581,-4.011) circle (0.20);
\fill[atomA, draw=atomLine, line width=0.35pt] (-6.701,-3.010) circle (0.20);
\fill[atomA, draw=atomLine, line width=0.35pt] (-6.820,-2.008) circle (0.20);
\fill[atomA, draw=atomLine, line width=0.35pt] (-6.940,-1.008) circle (0.20);
\fill[atomA, draw=atomLine, line width=0.35pt] (-6.060,-0.008) circle (0.20);
\fill[atomA, draw=atomLine, line width=0.35pt] (-6.180,0.992) circle (0.20);
\fill[atomA, draw=atomLine, line width=0.35pt] (-6.299,1.990) circle (0.20);
\fill[atomA, draw=atomLine, line width=0.35pt] (-6.419,2.989) circle (0.20);
\fill[atomA, draw=atomLine, line width=0.35pt] (-5.539,3.987) circle (0.20);
\fill[atomA, draw=atomLine, line width=0.35pt] (-5.658,4.984) circle (0.20);
\fill[atomA, draw=atomLine, line width=0.35pt] (-5.778,5.980) circle (0.20);
\fill[atomA, draw=atomLine, line width=0.35pt] (-5.897,6.976) circle (0.20);
\fill[atomA, draw=atomLine, line width=0.35pt] (-6.216,-6.922) circle (0.20);
\fill[atomA, draw=atomLine, line width=0.35pt] (-6.337,-5.919) circle (0.20);
\fill[atomA, draw=atomLine, line width=0.35pt] (-6.457,-4.916) circle (0.20);
\fill[atomA, draw=atomLine, line width=0.35pt] (-5.578,-3.913) circle (0.20);
\fill[atomA, draw=atomLine, line width=0.35pt] (-5.699,-2.911) circle (0.20);
\fill[atomA, draw=atomLine, line width=0.35pt] (-5.819,-1.910) circle (0.20);
\fill[atomA, draw=atomLine, line width=0.35pt] (-5.940,-0.910) circle (0.20);
\fill[atomA, draw=atomLine, line width=0.35pt] (-5.060,0.090) circle (0.20);
\fill[atomA, draw=atomLine, line width=0.35pt] (-5.181,1.090) circle (0.20);
\fill[atomA, draw=atomLine, line width=0.35pt] (-5.301,2.089) circle (0.20);
\fill[atomA, draw=atomLine, line width=0.35pt] (-5.422,3.087) circle (0.20);
\fill[atomA, draw=atomLine, line width=0.35pt] (-4.543,4.084) circle (0.20);
\fill[atomA, draw=atomLine, line width=0.35pt] (-4.663,5.081) circle (0.20);
\fill[atomA, draw=atomLine, line width=0.35pt] (-4.784,6.078) circle (0.20);
\fill[atomA, draw=atomLine, line width=0.35pt] (-4.904,7.073) circle (0.20);
\fill[atomA, draw=atomLine, line width=0.35pt] (-5.211,-6.820) circle (0.20);
\fill[atomA, draw=atomLine, line width=0.35pt] (-5.331,-5.816) circle (0.20);
\fill[atomA, draw=atomLine, line width=0.35pt] (-5.453,-4.814) circle (0.20);
\fill[atomA, draw=atomLine, line width=0.35pt] (-4.575,-3.811) circle (0.20);
\fill[atomA, draw=atomLine, line width=0.35pt] (-4.697,-2.808) circle (0.20);
\fill[atomA, draw=atomLine, line width=0.35pt] (-4.817,-1.807) circle (0.20);
\fill[atomA, draw=atomLine, line width=0.35pt] (-4.939,-0.807) circle (0.20);
\fill[atomA, draw=atomLine, line width=0.35pt] (-4.061,0.193) circle (0.20);
\fill[atomA, draw=atomLine, line width=0.35pt] (-4.183,1.193) circle (0.20);
\fill[atomA, draw=atomLine, line width=0.35pt] (-4.303,2.192) circle (0.20);
\fill[atomA, draw=atomLine, line width=0.35pt] (-4.425,3.189) circle (0.20);
\fill[atomA, draw=atomLine, line width=0.35pt] (-3.547,4.186) circle (0.20);
\fill[atomA, draw=atomLine, line width=0.35pt] (-3.669,5.184) circle (0.20);
\fill[atomA, draw=atomLine, line width=0.35pt] (-3.789,6.180) circle (0.20);
\fill[atomA, draw=atomLine, line width=0.35pt] (-3.910,7.174) circle (0.20);
\fill[atomA, draw=atomLine, line width=0.35pt] (-4.208,-6.713) circle (0.20);
\fill[atomA, draw=atomLine, line width=0.35pt] (-4.325,-5.709) circle (0.20);
\fill[atomA, draw=atomLine, line width=0.35pt] (-4.447,-4.708) circle (0.20);
\fill[atomA, draw=atomLine, line width=0.35pt] (-3.574,-3.706) circle (0.20);
\fill[atomA, draw=atomLine, line width=0.35pt] (-3.697,-2.700) circle (0.20);
\fill[atomA, draw=atomLine, line width=0.35pt] (-3.814,-1.699) circle (0.20);
\fill[atomA, draw=atomLine, line width=0.35pt] (-3.937,-0.701) circle (0.20);
\fill[atomA, draw=atomLine, line width=0.35pt] (-3.063,0.299) circle (0.20);
\fill[atomA, draw=atomLine, line width=0.35pt] (-3.186,1.301) circle (0.20);
\fill[atomA, draw=atomLine, line width=0.35pt] (-3.303,2.300) circle (0.20);
\fill[atomA, draw=atomLine, line width=0.35pt] (-3.426,3.294) circle (0.20);
\fill[atomA, draw=atomLine, line width=0.35pt] (-2.553,4.292) circle (0.20);
\fill[atomA, draw=atomLine, line width=0.35pt] (-2.675,5.291) circle (0.20);
\fill[atomA, draw=atomLine, line width=0.35pt] (-2.792,6.287) circle (0.20);
\fill[atomA, draw=atomLine, line width=0.35pt] (-2.915,7.279) circle (0.20);
\fill[atomA, draw=atomLine, line width=0.35pt] (-3.208,-6.598) circle (0.20);
\fill[atomA, draw=atomLine, line width=0.35pt] (-3.315,-5.594) circle (0.20);
\fill[atomA, draw=atomLine, line width=0.35pt] (-3.438,-4.601) circle (0.20);
\fill[atomA, draw=atomLine, line width=0.35pt] (-2.578,-3.598) circle (0.20);
\fill[atomA, draw=atomLine, line width=0.35pt] (-2.701,-2.585) circle (0.20);
\fill[atomA, draw=atomLine, line width=0.35pt] (-2.807,-1.583) circle (0.20);
\fill[atomA, draw=atomLine, line width=0.35pt] (-2.930,-0.594) circle (0.20);
\fill[atomA, draw=atomLine, line width=0.35pt] (-2.070,0.406) circle (0.20);
\fill[atomA, draw=atomLine, line width=0.35pt] (-2.193,1.417) circle (0.20);
\fill[atomA, draw=atomLine, line width=0.35pt] (-2.299,2.415) circle (0.20);
\fill[atomA, draw=atomLine, line width=0.35pt] (-2.422,3.402) circle (0.20);
\fill[atomA, draw=atomLine, line width=0.35pt] (-1.562,4.399) circle (0.20);
\fill[atomA, draw=atomLine, line width=0.35pt] (-1.685,5.406) circle (0.20);
\fill[atomA, draw=atomLine, line width=0.35pt] (-1.792,6.402) circle (0.20);
\fill[atomA, draw=atomLine, line width=0.35pt] (-2.221,-6.473) circle (0.20);
\fill[atomA, draw=atomLine, line width=0.35pt] (-2.293,-5.469) circle (0.20);
\fill[atomA, draw=atomLine, line width=0.35pt] (-2.411,-4.495) circle (0.20);
\fill[atomA, draw=atomLine, line width=0.35pt] (-1.598,-3.492) circle (0.20);
\fill[atomA, draw=atomLine, line width=0.35pt] (-1.716,-2.460) circle (0.20);
\fill[atomA, draw=atomLine, line width=0.35pt] (-1.789,-1.458) circle (0.20);
\fill[atomA, draw=atomLine, line width=0.35pt] (-1.906,-0.487) circle (0.20);
\fill[atomA, draw=atomLine, line width=0.35pt] (-1.094,0.513) circle (0.20);
\fill[atomA, draw=atomLine, line width=0.35pt] (-1.211,1.542) circle (0.20);
\fill[atomA, draw=atomLine, line width=0.35pt] (-1.284,2.540) circle (0.20);
\fill[atomA, draw=atomLine, line width=0.35pt] (-1.402,3.508) circle (0.20);
\fill[atomA, draw=atomLine, line width=0.35pt] (-0.589,4.505) circle (0.20);
\fill[atomA, draw=atomLine, line width=0.35pt] (-0.707,5.531) circle (0.20);
\fill[atomA, draw=atomLine, line width=0.35pt] (-0.779,6.527) circle (0.20);
\fill[atomA, draw=atomLine, line width=0.35pt] (-1.244,-6.356) circle (0.20);
\fill[atomA, draw=atomLine, line width=0.35pt] (-1.261,-5.352) circle (0.20);
\fill[atomA, draw=atomLine, line width=0.35pt] (-1.317,-4.365) circle (0.20);
\fill[atomA, draw=atomLine, line width=0.35pt] (-0.686,-3.362) circle (0.20);
\fill[atomA, draw=atomLine, line width=0.35pt] (-0.742,-2.343) circle (0.20);
\fill[atomA, draw=atomLine, line width=0.35pt] (-0.759,-1.341) circle (0.20);
\fill[atomA, draw=atomLine, line width=0.35pt] (-0.816,-0.357) circle (0.20);
\fill[atomA, draw=atomLine, line width=0.35pt] (-0.184,0.643) circle (0.20);
\fill[atomA, draw=atomLine, line width=0.35pt] (-0.241,1.659) circle (0.20);
\fill[atomA, draw=atomLine, line width=0.35pt] (-0.258,2.657) circle (0.20);
\fill[atomA, draw=atomLine, line width=0.35pt] (-0.314,3.638) circle (0.20);
\fill[atomA, draw=atomLine, line width=0.35pt] (0.317,4.635) circle (0.20);
\fill[atomA, draw=atomLine, line width=0.35pt] (0.261,5.648) circle (0.20);
\fill[atomA, draw=atomLine, line width=0.35pt] (0.244,6.644) circle (0.20);
\fill[atomA, draw=atomLine, line width=0.35pt] (-0.256,-6.356) circle (0.20);
\fill[atomA, draw=atomLine, line width=0.35pt] (-0.239,-5.352) circle (0.20);
\fill[atomA, draw=atomLine, line width=0.35pt] (-0.183,-4.365) circle (0.20);
\fill[atomA, draw=atomLine, line width=0.35pt] (0.186,-3.362) circle (0.20);
\fill[atomA, draw=atomLine, line width=0.35pt] (0.242,-2.343) circle (0.20);
\fill[atomA, draw=atomLine, line width=0.35pt] (0.259,-1.341) circle (0.20);
\fill[atomA, draw=atomLine, line width=0.35pt] (0.316,-0.357) circle (0.20);
\fill[atomA, draw=atomLine, line width=0.35pt] (0.684,0.643) circle (0.20);
\fill[atomA, draw=atomLine, line width=0.35pt] (0.741,1.659) circle (0.20);
\fill[atomA, draw=atomLine, line width=0.35pt] (0.758,2.657) circle (0.20);
\fill[atomA, draw=atomLine, line width=0.35pt] (0.814,3.638) circle (0.20);
\fill[atomA, draw=atomLine, line width=0.35pt] (1.183,4.635) circle (0.20);
\fill[atomA, draw=atomLine, line width=0.35pt] (1.239,5.648) circle (0.20);
\fill[atomA, draw=atomLine, line width=0.35pt] (1.256,6.644) circle (0.20);
\fill[atomA, draw=atomLine, line width=0.35pt] (0.721,-6.473) circle (0.20);
\fill[atomA, draw=atomLine, line width=0.35pt] (0.793,-5.469) circle (0.20);
\fill[atomA, draw=atomLine, line width=0.35pt] (0.911,-4.495) circle (0.20);
\fill[atomA, draw=atomLine, line width=0.35pt] (1.098,-3.492) circle (0.20);
\fill[atomA, draw=atomLine, line width=0.35pt] (1.216,-2.460) circle (0.20);
\fill[atomA, draw=atomLine, line width=0.35pt] (1.289,-1.458) circle (0.20);
\fill[atomA, draw=atomLine, line width=0.35pt] (1.406,-0.487) circle (0.20);
\fill[atomA, draw=atomLine, line width=0.35pt] (1.594,0.513) circle (0.20);
\fill[atomA, draw=atomLine, line width=0.35pt] (1.711,1.542) circle (0.20);
\fill[atomA, draw=atomLine, line width=0.35pt] (1.784,2.540) circle (0.20);
\fill[atomA, draw=atomLine, line width=0.35pt] (1.902,3.508) circle (0.20);
\fill[atomA, draw=atomLine, line width=0.35pt] (2.089,4.505) circle (0.20);
\fill[atomA, draw=atomLine, line width=0.35pt] (2.207,5.531) circle (0.20);
\fill[atomA, draw=atomLine, line width=0.35pt] (2.279,6.527) circle (0.20);
\fill[atomA, draw=atomLine, line width=0.35pt] (1.708,-6.598) circle (0.20);
\fill[atomA, draw=atomLine, line width=0.35pt] (1.815,-5.594) circle (0.20);
\fill[atomA, draw=atomLine, line width=0.35pt] (1.938,-4.601) circle (0.20);
\fill[atomA, draw=atomLine, line width=0.35pt] (2.078,-3.598) circle (0.20);
\fill[atomA, draw=atomLine, line width=0.35pt] (2.201,-2.585) circle (0.20);
\fill[atomA, draw=atomLine, line width=0.35pt] (2.307,-1.583) circle (0.20);
\fill[atomA, draw=atomLine, line width=0.35pt] (2.430,-0.594) circle (0.20);
\fill[atomA, draw=atomLine, line width=0.35pt] (2.570,0.406) circle (0.20);
\fill[atomA, draw=atomLine, line width=0.35pt] (2.693,1.417) circle (0.20);
\fill[atomA, draw=atomLine, line width=0.35pt] (2.799,2.415) circle (0.20);
\fill[atomA, draw=atomLine, line width=0.35pt] (2.922,3.402) circle (0.20);
\fill[atomA, draw=atomLine, line width=0.35pt] (3.062,4.399) circle (0.20);
\fill[atomA, draw=atomLine, line width=0.35pt] (3.185,5.406) circle (0.20);
\fill[atomA, draw=atomLine, line width=0.35pt] (3.292,6.402) circle (0.20);
\fill[atomA, draw=atomLine, line width=0.35pt] (2.708,-6.713) circle (0.20);
\fill[atomA, draw=atomLine, line width=0.35pt] (2.825,-5.709) circle (0.20);
\fill[atomA, draw=atomLine, line width=0.35pt] (2.947,-4.708) circle (0.20);
\fill[atomA, draw=atomLine, line width=0.35pt] (3.074,-3.706) circle (0.20);
\fill[atomA, draw=atomLine, line width=0.35pt] (3.197,-2.700) circle (0.20);
\fill[atomA, draw=atomLine, line width=0.35pt] (3.314,-1.699) circle (0.20);
\fill[atomA, draw=atomLine, line width=0.35pt] (3.437,-0.701) circle (0.20);
\fill[atomA, draw=atomLine, line width=0.35pt] (3.563,0.299) circle (0.20);
\fill[atomA, draw=atomLine, line width=0.35pt] (3.686,1.301) circle (0.20);
\fill[atomA, draw=atomLine, line width=0.35pt] (3.803,2.300) circle (0.20);
\fill[atomA, draw=atomLine, line width=0.35pt] (3.926,3.294) circle (0.20);
\fill[atomA, draw=atomLine, line width=0.35pt] (4.053,4.292) circle (0.20);
\fill[atomA, draw=atomLine, line width=0.35pt] (4.175,5.291) circle (0.20);
\fill[atomA, draw=atomLine, line width=0.35pt] (4.292,6.287) circle (0.20);
\fill[atomA, draw=atomLine, line width=0.35pt] (4.415,7.279) circle (0.20);
\fill[atomA, draw=atomLine, line width=0.35pt] (3.711,-6.820) circle (0.20);
\fill[atomA, draw=atomLine, line width=0.35pt] (3.831,-5.816) circle (0.20);
\fill[atomA, draw=atomLine, line width=0.35pt] (3.953,-4.814) circle (0.20);
\fill[atomA, draw=atomLine, line width=0.35pt] (4.075,-3.811) circle (0.20);
\fill[atomA, draw=atomLine, line width=0.35pt] (4.197,-2.808) circle (0.20);
\fill[atomA, draw=atomLine, line width=0.35pt] (4.317,-1.807) circle (0.20);
\fill[atomA, draw=atomLine, line width=0.35pt] (4.439,-0.807) circle (0.20);
\fill[atomA, draw=atomLine, line width=0.35pt] (4.561,0.193) circle (0.20);
\fill[atomA, draw=atomLine, line width=0.35pt] (4.683,1.193) circle (0.20);
\fill[atomA, draw=atomLine, line width=0.35pt] (4.803,2.192) circle (0.20);
\fill[atomA, draw=atomLine, line width=0.35pt] (4.925,3.189) circle (0.20);
\fill[atomA, draw=atomLine, line width=0.35pt] (5.047,4.186) circle (0.20);
\fill[atomA, draw=atomLine, line width=0.35pt] (5.169,5.184) circle (0.20);
\fill[atomA, draw=atomLine, line width=0.35pt] (5.289,6.180) circle (0.20);
\fill[atomA, draw=atomLine, line width=0.35pt] (5.410,7.174) circle (0.20);
\fill[atomA, draw=atomLine, line width=0.35pt] (4.716,-6.922) circle (0.20);
\fill[atomA, draw=atomLine, line width=0.35pt] (4.837,-5.919) circle (0.20);
\fill[atomA, draw=atomLine, line width=0.35pt] (4.957,-4.916) circle (0.20);
\fill[atomA, draw=atomLine, line width=0.35pt] (5.078,-3.913) circle (0.20);
\fill[atomA, draw=atomLine, line width=0.35pt] (5.199,-2.911) circle (0.20);
\fill[atomA, draw=atomLine, line width=0.35pt] (5.319,-1.910) circle (0.20);
\fill[atomA, draw=atomLine, line width=0.35pt] (5.440,-0.910) circle (0.20);
\fill[atomA, draw=atomLine, line width=0.35pt] (5.560,0.090) circle (0.20);
\fill[atomA, draw=atomLine, line width=0.35pt] (5.681,1.090) circle (0.20);
\fill[atomA, draw=atomLine, line width=0.35pt] (5.801,2.089) circle (0.20);
\fill[atomA, draw=atomLine, line width=0.35pt] (5.922,3.087) circle (0.20);
\fill[atomA, draw=atomLine, line width=0.35pt] (6.043,4.084) circle (0.20);
\fill[atomA, draw=atomLine, line width=0.35pt] (6.163,5.081) circle (0.20);
\fill[atomA, draw=atomLine, line width=0.35pt] (6.284,6.078) circle (0.20);
\fill[atomA, draw=atomLine, line width=0.35pt] (6.404,7.073) circle (0.20);
\fill[atomA, draw=atomLine, line width=0.35pt] (5.722,-7.020) circle (0.20);
\fill[atomA, draw=atomLine, line width=0.35pt] (5.842,-6.016) circle (0.20);
\fill[atomA, draw=atomLine, line width=0.35pt] (5.961,-5.013) circle (0.20);
\fill[atomA, draw=atomLine, line width=0.35pt] (6.081,-4.011) circle (0.20);
\fill[atomA, draw=atomLine, line width=0.35pt] (6.201,-3.010) circle (0.20);
\fill[atomA, draw=atomLine, line width=0.35pt] (6.320,-2.008) circle (0.20);
\fill[atomA, draw=atomLine, line width=0.35pt] (6.440,-1.008) circle (0.20);
\fill[atomA, draw=atomLine, line width=0.35pt] (6.560,-0.008) circle (0.20);
\fill[atomA, draw=atomLine, line width=0.35pt] (6.680,0.992) circle (0.20);
\fill[atomA, draw=atomLine, line width=0.35pt] (6.799,1.990) circle (0.20);
\fill[atomA, draw=atomLine, line width=0.35pt] (6.919,2.989) circle (0.20);
\fill[atomA, draw=atomLine, line width=0.35pt] (7.039,3.987) circle (0.20);
\fill[atomA, draw=atomLine, line width=0.35pt] (7.158,4.984) circle (0.20);
\fill[atomA, draw=atomLine, line width=0.35pt] (7.278,5.980) circle (0.20);
\fill[atomA, draw=atomLine, line width=0.35pt] (6.728,-7.113) circle (0.20);
\fill[atomA, draw=atomLine, line width=0.35pt] (6.847,-6.110) circle (0.20);
\fill[atomA, draw=atomLine, line width=0.35pt] (6.965,-5.107) circle (0.20);
\fill[atomA, draw=atomLine, line width=0.35pt] (7.084,-4.105) circle (0.20);
\fill[atomA, draw=atomLine, line width=0.35pt] (7.203,-3.104) circle (0.20);
\end{scope}
\dislT[fmRed]{(0.000,-3.055)}{0}
\dislT[fmRed]{(0.000,0.945)}{0}
\dislT[fmRed]{(0.000,4.945)}{0}
\coordinate (TB0) at (0,0.995);
\coordinate (TB1) at (0,4.995);

  \draw[dimline, fmRed] ($(TB0)+(1.3,0)$) -- node[right, slbl, fmRed]{$h$} ($(TB1)+(1.3,0)$);
\end{tikzpicture}
\caption{Малоугловая граница наклона~— вертикальная стенка краевых
дислокаций (значки $\perp$), расположенных через $h=4b$. Атомные ряды по
разные стороны стенки повёрнуты на угол $\varphi=b/h\approx14^\circ$.
Положения атомов рассчитаны суперпозицией упругих полей дислокаций стенки}
\label{fig:tilt-wall}
\end{figure}

\begin{opredelenie}
Слабо разориентированные области внутри зерна, разделённые
дислокационными стенками, называют \term{субзёрнами} (блоками,
фрагментами), а сами стенки~— \term{субграницами} (малоугловыми
границами).
\end{opredelenie}

При небольших напряжениях $h>50b$, и $\varphi<1/50$~рад~$\approx1^\circ$.
С ростом напряжений плотность дислокаций в стенках растёт, $h$
уменьшается, и $\varphi$ может достигать $5$--$10^\circ$~— субграница
превращается в границу зерна. Так при деформации зёрна дробятся.

Стенка из \emph{винтовых} дислокаций (точнее, сетка из двух семейств
винтовых дислокаций) образует \term{границу кручения}: соседние части
кристалла повёрнуты на угол $\varphi=b/h$ вокруг оси, перпендикулярной
границе. Границы наклона и кручения~— два крайних случая устойчивых
границ, не создающих напряжений на больших расстояниях; на практике часто
встречаются границы смешанного типа, напряжения около которых
релаксированы не полностью.

\subsubsection{Конечные стенки и полигонизация}

Бесконечную стенку можно представить в монокристалле; в поликристалле
стенки конечны. У конечной стенки из $n$ дислокаций высотой $H=nh$
полной компенсации полей нет. Её поле удобно описать на трёх масштабах:
\begin{itemize}
\item у самой оси дислокации ($x,y\ll h$) работает только её собственное
поле;
\item вблизи стенки и далеко от её концов ($x<h$, $y\ll H$) поле почти
такое же, как у бесконечной стенки~— периодическое и быстро затухающее;
\item далеко от стенки ($x\gg H$) она выглядит как одна дислокация с
вектором $n\vb b$, и её поле убывает как $1/r$~— оно \emph{дальнодействующее}.
\end{itemize}

Дальнодействующее поле конечной стенки притягивает одноимённые
дислокации из окружающего объёма в её плоскость (и выталкивает
разноимённые). Чем больше дислокаций в стенке, тем сильнее поле~— стенка
растёт \emph{автокаталитически}, с самоускорением.

\begin{opredelenie}
\term{Полигонизация}~— самопроизвольная перестройка дислокаций в более
устойчивые конфигурации (стенки, субграницы), не сопровождающаяся
образованием или движением границ зёрен. Движущая сила~— упругая энергия
дислокаций; скорость обычно определяется самым медленным процессом~—
переползанием.
\end{opredelenie}

Классический пример~— отжиг изогнутого кристалла
(рис.~\ref{fig:polygonization}). При изгибе в нём остаётся избыток
краевых дислокаций одного знака, распределённых по плоскостям
скольжения. При нагреве они переползают и выстраиваются в стенки;
изогнутый кристалл превращается в «многоугольник» из слабо
разориентированных субзёрен с плоскими участками. Полигонизацию наблюдают
при нагреве после деформации, при горячей деформации и ползучести, при
охлаждении после затвердевания.

\begin{figure}[ht]
\centering
\begin{tikzpicture}[scale=1.15]
  \begin{scope}
    \foreach \r in {3.0,3.4,3.8,4.2}{\draw[fmBlue!60] (60:\r) arc (60:120:\r);}
    \foreach \ang/\r in {68/3.2,75/3.6,83/3.25,90/3.95,97/3.55,105/3.2,112/3.75,79/4.0,101/4.0,86/3.6,108/3.6,72/3.95}{
      \dislT[fmRed]{(\ang:\r)}{\ang-90}}
    \node[lbl] at (90:2.4) {а};
  \end{scope}
  \begin{scope}[shift={(6.2,0)}]
    \foreach \r in {3.0,3.4,3.8,4.2}{
      \draw[fmBlue!60] (60:\r) -- (75:\r) -- (90:\r) -- (105:\r) -- (120:\r);}
    \foreach \ang in {75,90,105}{
      \foreach \r in {3.2,3.6,4.0}{\dislT[fmRed]{(\ang:\r)}{\ang-90}}}
    \node[lbl] at (90:2.4) {б};
  \end{scope}
\end{tikzpicture}
\caption{Полигонизация: \textit{а}~— изогнутый кристалл с избыточными
краевыми дислокациями одного знака; \textit{б}~— после отжига дислокации
выстроились в стенки (субграницы), разделяющие плоские участки~— субзёрна}
\label{fig:polygonization}
\end{figure}

\begin{podrobnee}[Системы субграниц и дисклинации]
Чтобы описать систему субграниц, не вдаваясь в строение каждой,
субграницу рассматривают как поверхность, создающую разворот на угол
$\omega$ и обладающую удельной энергией $\sigma(\omega)$. Отдельные
площадки субграниц стыкуются попарно, по три (тройной стык) или
обрываются внутри решётки.

\emph{Дисклинации.} Край обрывающейся субграницы~— особый линейный
дефект. Если обойти его контуром, то разрыва смещений нет ($\oint\d\vb
u=0$), но накапливается \emph{поворот} на угол $\omega$:
$\oint\d\vb\omega=\vb\omega$. Такой дефект называют \term{дисклинацией}.
Её характеризуют направление линии $\vb l$ и \emph{вектор Франка}
$\vb\omega$ (вектор поворота). Для дисклинаций, как и для дислокаций,
действуют законы сохранения: дисклинация непрерывна (внутри кристалла
образует петлю), $\vb\omega$ постоянен вдоль неё, при реакциях
$\vb\omega_3=\vb\omega_1+\vb\omega_2$, в узлах $\sum\vb\omega_i=0$. В
отличие от дислокаций, набор $\vb\omega$ не дискретен: угол плавно
регулируется плотностью дислокаций в субгранице. Типы дисклинаций:
\emph{клиновые} ($\vb\omega\parallel\vb l$) и два типа \emph{дисклинаций
кручения} ($\vb\omega\parallel\vb n$~— нормали к плоскости петли, и
$\vb\omega\parallel\vb l\times\vb n$). Петля, очерчивающая площадку
границы наклона, состоит из двух клиновых участков и двух участков
кручения; петля, очерчивающая границу кручения,~— только из дисклинаций
кручения.

Энергия поля дисклинации квадратично зависит от $\omega$, поэтому
дальнодействующих полей нет только у тех стыков субграниц, где
$\sum\vb\omega_i=0$. Пример: кубик, очерченный одной стопкой
дислокационных петель, имеет четыре субграницы и открытые верх и низ; на
его рёбрах $\sum\omega_i\neq0$, и есть дальнее поле. Если закрыть
открытые грани второй системой петель, на всех рёбрах $\sum\omega_i=0$,
и дальнодействие исчезает.

\emph{Равновесие стыка.} Если энергию субграницы трактовать как
поверхностное натяжение, то в тройном стыке, как для плёнок,
$\vb\sigma_1+\vb\sigma_2+\vb\sigma_3=0$. «Устойчивость в большом»
($\sum\vb\omega_i=0$) устанавливается перемещениями дислокаций на
большие расстояния, «устойчивость в малом» ($\sum\vb\sigma_i=0$)~—
местными движениями границ около стыка. Слияние двух субграниц всегда
даёт выигрыш энергии (большой, если оси разворота разные, и малый, если
одинаковые), поэтому никакая система субграниц не устойчива абсолютно
по отношению к укрупнению субзёрен.
\end{podrobnee}

\begin{itogi}
\begin{itemize}
\item Скопление: неустойчиво, концентрирует напряжение в $N$ раз;
$d=2Ab/N\tau$, $L\approx2NAb/\tau$. Отсюда формула Холла--Петча
$\tau_\text{т}=\tau_0+kl^{-1/2}$: мелкое зерно~— высокая прочность.
\item Стенка: устойчива, поля гасятся на расстоянии $\sim h$; энергия
меньше в 3--4 раза. Стенка краевых дислокаций~— граница наклона
($\varphi=b/h$), винтовых~— граница кручения.
\item Стенки образуются переползанием, поэтому при $T\gtrsim(0{,}3\text{--}0{,}5)\Tpl$.
Полигонизация~— перестройка дислокаций в субграницы.
\end{itemize}
\end{itogi}

\begin{proverka}
\item Почему скопление не может образоваться перегруппировкой
дислокаций?
\item Во сколько раз скопление из 50 дислокаций усиливает приложенное
напряжение?
\item Объясните, почему мелкозернистый металл прочнее крупнозернистого.
\item Какова разориентация субзёрен, если расстояние между дислокациями в
стенке $h=100b$?
\item Почему полигонизация не идёт при низких температурах?
\end{proverka}
